\clearpage
\section[HKO Local Maximality]{HKO Local Maximality Among Ten-Facet Polytopes}
\label{sec:hko-local-maximum}

Let \(\mathcal K\) be the space of convex bodies in \(\mathbb R^4\), equipped
with the Hausdorff metric, and let \(\mathcal P_{10}\subset\mathcal K\) be the
subspace of convex polytopes with exactly ten facets. Let
\(K_{\mathrm{HKO}}\) be the HKO polytope from
\cite[Proposition~1.3]{HaimKislevOstrover2024}: the Lagrangian product of the
two regular pentagons with circumradius \(1\) used there. We use the centered
representative, so \(0\in\operatorname{int}(K_{\mathrm{HKO}})\). It lies in
\(\mathcal P_{10}\), and the computation in
\cite[Proof of Theorem~1.2]{HaimKislevOstrover2024} gives
\[
  \operatorname{sys}(K_{\mathrm{HKO}})
  =
  \frac{\sqrt{5}+3}{5}
  >1.
\]
The result of this section is local: inside the ten-facet space, no
sufficiently small perturbation improves the systolic ratio once the natural
symmetries have been removed.

\begin{remark}[Certificate status]
The exact SageMath verifier checks the finite certificate described below. The
remaining review question is whether the Hausdorff-to-chart reduction and the
quotient-slice translation are accepted as final theorem-strength thesis text.
\end{remark}

\begin{theorem}[HKO is locally maximal in the ten-facet space]
\label{thm:hko-ten-facet-local-maximum}
There is a Hausdorff neighborhood \(\mathcal U\) of \(K_{\mathrm{HKO}}\) in
\(\mathcal P_{10}\) such that
\[
  \operatorname{sys}(K)\leq \operatorname{sys}(K_{\mathrm{HKO}})
  \qquad (K\in\mathcal U).
\]
Moreover, equality can occur near \(K_{\mathrm{HKO}}\) only along the natural
symmetries: if \(K\in\mathcal U\) has
\[
  \operatorname{sys}(K)=\operatorname{sys}(K_{\mathrm{HKO}}),
\]
then \(K\) lies in the local orbit of \(K_{\mathrm{HKO}}\) under translations,
positive scaling, and linear symplectic transformations.
\end{theorem}

\clearpage
\subsection{Reduction to the certificate}
\label{subsec:hko-local-maximum-decision-problem}

The proof is local, and the first step is not computational. The following
elementary reduction is the bridge from the Hausdorff theorem statement to the
dual-vertex coordinates used by the certificate.

\begin{lemma}[Hausdorff-local ten-facet chart]
\label{lem:hko-hausdorff-local-ten-facet-chart}
After shrinking the Hausdorff neighborhood of \(K_{\mathrm{HKO}}\) in
\(\mathcal P_{10}\), every nearby polytope has facets labelled by the ten
facets of \(K_{\mathrm{HKO}}\). With these labels it has a unique
dual-vertex representation
\[
  K(a)=
  \{x\in\mathbb R^4:\langle a_i,x\rangle\le 1,\ i=1,\ldots,10\},
  \qquad
  a=(a_1,\ldots,a_{10})\in(\mathbb R^4)^{10}.
\]
In these coordinates the supporting inequalities stay facet-defining,
boundedness and nonempty interior persist, and the HKO face combinatorics is
unchanged. Conversely, a sufficiently small coordinate neighborhood of the HKO
dual vertices parametrizes this labelled Hausdorff-local stratum.
\end{lemma}

\begin{proof}
The centered HKO representative contains the origin in its interior. After
shrinking the Hausdorff neighborhood, the origin remains in the interior of
every nearby polytope, and polarity is continuous on this neighborhood. Thus a
polytope in \(\mathcal P_{10}\) has ten facets if and only if its polar has ten
vertices.

Since \(K_{\mathrm{HKO}}\) is simple, its irredundant inequalities are stable
under small labelled perturbations. Equivalently,
\(K_{\mathrm{HKO}}^\circ\) is a simplicial polytope with ten vertices. For each
vertex of \(K_{\mathrm{HKO}}^\circ\), choose a linear functional in the interior
of its normal cone. If a ten-vertex polytope \(L\) is sufficiently
Hausdorff-close to \(K_{\mathrm{HKO}}^\circ\), then the maximizers of these ten
functionals lie in disjoint small neighborhoods of the corresponding vertices.
Since \(L\) has only ten vertices, this labels all vertices of \(L\). Applying
this to \(L=K^\circ\), every Hausdorff-near \(K\in\mathcal P_{10}\) has ten
facets corresponding to the ten facets of \(K_{\mathrm{HKO}}\), and the
corresponding polar vertices give the vectors \(a_i\).

It remains only to shrink to the stable simple-polytope chart. We use the
standard local stability of simple labelled \(H\)-representations: after
shrinking the labelled coefficients, a simple polytope keeps the same face
lattice. In the present chart this is witnessed by finitely many open
conditions. At \(K_{\mathrm{HKO}}\), each vertex is the transverse intersection
of its four incident supporting hyperplanes, every nonincident inequality has
strict slack at that vertex, and the labelled inequalities are irredundant.
These determinant, slack, irredundancy, boundedness, and origin-interior
conditions remain true for all sufficiently small labelled perturbations, so the
HKO face combinatorics is unchanged. Conversely, small perturbations of the
labelled polar vertices keeping these same open conditions give exactly the
labelled local stratum after shrinking once more.
\end{proof}

This local stability statement is separate from the SageMath certificate: the
certificate begins only after the Hausdorff-local problem has been put in these
labelled dual-vertex coordinates. By
Lemma~\ref{lem:hko-hausdorff-local-ten-facet-chart}, a coordinate-local
conclusion in the chart proves the stated local result in \(\mathcal P_{10}\).

The natural symmetries of the systolic ratio induce a local action on the
dual-vertex coordinates. A translation by \(y\in\mathbb R^4\) sends
\[
  a_i\mapsto \frac{a_i}{1+\langle a_i,y\rangle},
\]
positive scaling by \(\rho>0\) sends \(a_i\mapsto \rho^{-1}a_i\), and a
linear symplectic map \(M\) sends \(a_i\mapsto M^{-T}a_i\). Differentiating
this induced action at \(K_{\mathrm{HKO}}\) gives a \(15\)-dimensional
symmetry tangent space \(T_{\mathrm{orb}}\subset(\mathbb R^4)^{10}\).

Choose a linear complement \(S\) to \(T_{\mathrm{orb}}\). The exact verifier
checks that \(\dim T_{\mathrm{orb}}=15\), so \(\dim S=25\). After shrinking the
neighborhoods, the map
\[
  (g,s)\longmapsto g\cdot(a_{\mathrm{HKO}}+s)
\]
from a neighborhood of the identity in
\(\mathbb R^4\rtimes(\mathbb R_{>0}\times Sp(4,\mathbb R)^0)\) and a
neighborhood of \(0\) in \(S\) is a local diffeomorphism onto a neighborhood of
\(a_{\mathrm{HKO}}\) in the chart. After shrinking, each nearby point has a
unique representation \(g\cdot(a_{\mathrm{HKO}}+s)\) with \(g\) near the
identity and \(s\in S\) small. Thus it is enough to prove that
\(\operatorname{sys}\) has a strict local maximum on \(a_{\mathrm{HKO}}+S\). If
equality holds near HKO, the slice component must be \(s=0\), and the polytope
lies in the local symmetry orbit.

\subsection{Exact finite certificate}
\label{subsec:hko-local-maximum-computation-with-sagemath}

The exact certificate proves the slice statement by upper branches. As in the
finite-capacity algorithm from
Section~\ref{sec:quadratic-program-algorithm-hk2019}, a word \(\sigma\) fixes
a cyclic facet sequence, and the coefficients \(\beta\) are the feasible
weights that produce an admissible action for that sequence. For each selected
word \(\sigma\), the certificate contains exact data defining a feasible
section \(\beta_\sigma(a)\). The corresponding function
\[
  U_\sigma(a)
\]
is a smooth upper bound for \(\operatorname{sys}(K(a))\) near the HKO point,
and satisfies
\[
  U_\sigma(a_{\mathrm{HKO}})
  =
  \operatorname{sys}(K_{\mathrm{HKO}}).
\]
These feasible sections are not required to be optimizing branches for nearby
parameters. They are used only as upper branches touching the systolic ratio
at HKO.

This point is essential for the HKO certificate. The 26-row certificate uses
both nonsingular six-facet data and singular seven-facet data; the
nonsingular active rows alone do not span the \(25\)-dimensional quotient
slice. The proof therefore does not differentiate nearby KKT-critical
optimizers. Instead, for each selected word it fixes enough beta coordinates
so that the closure and normalization equations have an invertible minor. This
gives a smooth feasible beta section and hence a differentiable upper branch
even when the KKT system for an optimizing branch is singular.

The finite first-order criterion is the following. Let
\(r_\sigma=D U_\sigma(a_{\mathrm{HKO}})\). If the selected rows annihilate the
symmetry tangent space, have ambient row rank \(25\), and admit positive
coefficients \(\lambda_\sigma\) with
\[
  \sum_\sigma\lambda_\sigma=1,
  \qquad
  \sum_\sigma\lambda_\sigma r_\sigma=0,
\]
then every nonzero transverse direction has some upper branch with strictly
negative first derivative. Indeed, the ambient coordinate space has dimension
\(40\), and the symmetry tangent space has dimension \(15\), so rank \(25\)
rows that vanish on the symmetry tangent space span the cotangent space of any
transverse slice. Since the selected set is finite and the upper branches are
smooth, the first-order decrease has a uniform negative margin on the compact
unit sphere in the slice. Hence \(\operatorname{sys}\) decreases for every
sufficiently small nonzero transverse perturbation.

The finite certificate is a list of exact assertions over the ordered number
field determined by the real root \(0<t<1\) of
\[
  t^4-10t^2+5=0.
\]
It records the following proof-facing facts:
\[
\begin{array}{ll}
\text{ambient dimension} & 40\\
\text{selected feasible sections} & 26\\
\text{symmetry tangent rank} & 15\\
\text{upper-branch row rank} & 25\\
\text{convex certificate} &
  \lambda_\sigma>0,\quad
  \sum_\sigma\lambda_\sigma=1,\quad
  \sum_\sigma\lambda_\sigma r_\sigma=0.
\end{array}
\]
The certificate also checks positivity, closure and normalization of each beta
vector, invertibility of the selected feasible-section minor, equality of each
selected branch action with the HKO action, the feasible-section derivative
equation, the derivative row formula for \(U_\sigma\), and annihilation of the
symmetry tangent space. These checks enter the proof in three places. The
field, HKO geometry, capacity, volume, and volume derivative checks identify
the base point and the normalization used for \(\operatorname{sys}\). The
rowwise checks give the smooth feasible upper branches \(U_\sigma\) and
identify their derivatives at the HKO point. The global rank, annihilation, and
positive-convex-combination checks are exactly the finite first-order
certificate on the quotient slice.

Rust is used only to generate the finite witness file. SageMath reads that
witness, reconstructs the exact algebraic data in memory, and checks the
proof-facing equations. The full verifier packet is stored in
\path{experiments/hko-local-maximum/theorem/}. The mathematical implication
from these checked equations to
Theorem~\ref{thm:hko-ten-facet-local-maximum} is the finite upper-branch and
slice argument above; it is not a property of the Rust search procedure.

\begin{proof}[Proof of Theorem~\ref{thm:hko-ten-facet-local-maximum}]
Lemma~\ref{lem:hko-hausdorff-local-ten-facet-chart} reduces the Hausdorff-local
statement in \(\mathcal P_{10}\) to the labelled dual-vertex chart. The local
slice paragraph reduces the quotient-local statement to strict local
maximality on the transverse slice \(a_{\mathrm{HKO}}+S\). The exact finite
certificate gives that strict slice maximum by the upper-branch argument above.
Therefore every sufficiently small perturbation has systolic ratio at most
\(\operatorname{sys}(K_{\mathrm{HKO}})\). If equality holds, strictness on the
slice forces the slice coordinate to vanish, so the perturbation is in the
local orbit of \(K_{\mathrm{HKO}}\) under translations, positive scaling, and
linear symplectic transformations.
\end{proof}

\subsection{Empirical tests}
\label{subsec:hko-local-maximum-empirical-tests}

The proof of Theorem~\ref{thm:hko-ten-facet-local-maximum} uses the exact
feasible-section certificate, not numerical sampling. We also kept empirical
sanity checks around the HKO point: numerical first- and second-order
perturbation tests, and recorded \(F=11\) facet-addition/ascent tests. The
committed runs did not find a higher nearby value of \(\operatorname{sys}\).
These experiments support the broader local-stability picture, but they are not
substitutes for the exact certificate.
